\documentclass[10pt,a4paper]{article}

\usepackage[utf8]{inputenc}
\usepackage[T1]{fontenc}
\usepackage{lmodern}
\usepackage[a4paper,margin=1.6cm,top=1.6cm,bottom=1.8cm]{geometry}
\usepackage{amssymb}
\usepackage{amsmath}
\usepackage{xcolor}
\usepackage{tabularx}
\usepackage{array}
\usepackage{booktabs}
\usepackage{enumitem}
\usepackage{titlesec}
\usepackage{fancyhdr}
\usepackage{tcolorbox}
\usepackage[hidelinks]{hyperref}

% ---------- colours and tags ----------
\definecolor{slides}{RGB}{0,92,153}
\definecolor{book}{RGB}{170,60,0}
\definecolor{notes}{RGB}{40,120,60}
\definecolor{devkit}{RGB}{110,40,140}
\definecolor{rowgray}{gray}{0.95}

\newcommand{\cb}{\raisebox{-0.35ex}{\Large$\square$}}
\newcommand{\srctag}[2]{\colorbox{#1}{\textcolor{white}{\textbf{\scriptsize\,#2\,}}}}
\newcommand{\SLIDES}{\srctag{slides}{SLIDES}}
\newcommand{\BOOK}{\srctag{book}{BOOK}}
\newcommand{\NOTES}{\srctag{notes}{NOTES}}
\newcommand{\DEVKIT}{\srctag{devkit}{DEVKIT}}
\newcommand{\LAB}{\srctag{devkit}{LAB}}

\newcolumntype{C}{>{\centering\arraybackslash}p{0.6cm}}
\newcolumntype{T}{>{\centering\arraybackslash}p{1.45cm}}

% ---------- sections ----------
\titleformat{\section}{\large\bfseries\color{slides}}{}{0pt}{}[\vspace{-0.5ex}\titlerule]
\titlespacing*{\section}{0pt}{2.2ex}{1.2ex}
\titleformat{\subsection}{\normalsize\bfseries}{}{0pt}{}
\titlespacing*{\subsection}{0pt}{1.8ex}{0.8ex}

\setlength{\parindent}{0pt}
\setlength{\parskip}{0.5ex}
\renewcommand{\arraystretch}{1.25}
\setlist[itemize]{nosep,leftmargin=1.2em}

\pagestyle{fancy}
\fancyhf{}
\fancyhead[L]{\small ARob 2026/27 - Lab 3 reading plan}
\fancyhead[R]{\small Pedro}
\fancyfoot[C]{\small\thepage}
\renewcommand{\headrulewidth}{0.3pt}

\newtcolorbox{selfcheck}{colback=rowgray,colframe=gray!60,boxrule=0.4pt,
  left=4pt,right=4pt,top=3pt,bottom=3pt,title={\small\bfseries Can I do this without looking?},
  fonttitle=\small,coltitle=black,colbacktitle=gray!25}

\newtcolorbox{warnbox}{colback=orange!8,colframe=orange!70!black,boxrule=0.5pt,
  left=4pt,right=4pt,top=3pt,bottom=3pt,title={\small\bfseries Found in the materials (check before coding)},
  fonttitle=\small,coltitle=black,colbacktitle=orange!25}

\newcommand{\chk}[1]{\par\noindent\hangindent=1.6em\hangafter=1 \cb\hspace{0.6em}#1}

\begin{document}

{\LARGE\bfseries\color{slides} Lab 3 - Theory and DevKit reading plan}\\[0.6ex]
{\large Motion Control of the Parrot AR.Drone (MEAer, UAVs, 2026/27)}\\[0.4ex]
{\small Goal: arrive at the session knowing every concept behind guide Sections 2 and 3, and exactly what the two empty Simulink blocks expect. Tick each box when done.}

\subsection*{How to read this sheet}

\begin{tabularx}{\textwidth}{@{}l X@{}}
\SLIDES & A lecture deck PDF in \texttt{lectures-22-23/}. The page number is the \textbf{PDF page} of that deck (page 1 is the cover).\\[0.5ex]
\BOOK & \texttt{reference-book/uavbook.pdf} (Beard and McLain). Two page numbers: the one \textbf{printed} on the page, and the \textbf{PDF} page (PDF = printed + 8).\\[0.5ex]
\NOTES & A written summary PDF in \texttt{lectures-22-23/notes/} or \texttt{labs-25-26/notes/}. Not slides, not the book.\\[0.5ex]
\DEVKIT & A file of \texttt{Lab3/DevKit\_Lab3/}. The Simulink facts in Area 6 were read directly from the \texttt{.slx} files (they are zipped XML), not taken from the 25/26 notes, which never opened them.\\[0.5ex]
\LAB & Your own earlier lab material in \texttt{labs-26-27/Lab1/}.\\
\end{tabularx}

\subsection*{The files used}

\begin{tabularx}{\textwidth}{@{}l X r@{}}
\toprule
Short name & File (relative to the repository root) & Pages\\
\midrule
Guide & \texttt{labs-26-27/Lab3/26\_27\_UAVs\_Lab3.pdf} (same as 25/26 except dates) & 5\\
Ch2 deck & \texttt{lectures-22-23/22\_23\_UAVs\_Ch2\_Quadrotor\_Modeling.pdf} & 30\\
Ch3 deck & \texttt{lectures-22-23/22\_23\_UAVs\_Ch3\_Quadrotor\_Control\_Intro.pdf} & 44\\
Ch4 deck & \texttt{lectures-22-23/22\_23\_UAVs\_Ch4\_Modern\_Control\_Design.pdf} & 43\\
Ch7 deck & \texttt{lectures-22-23/22\_23\_UAVs\_Ch7\_Intro\_Nonlinear\_Control.pdf} (Lyapunov, LaSalle) & 46\\
Ch7-II deck & \texttt{lectures-22-23/22\_23\_UAVs\_Ch7\_Intro\_Nonlinear\_Control\_Part\_II.pdf} (adaptive, tracking) & 42\\
Ch8 deck & \texttt{lectures-22-23/22\_23\_UAVs\_Ch8\_Path\_Following.pdf} & 19\\
Book & \texttt{reference-book/uavbook.pdf} & 401\\
Notes & \texttt{lectures-22-23/notes/02-control.pdf}, \texttt{04-guidance-nonlinear.pdf}; \texttt{labs-25-26/notes/03-lab3.pdf} & 9, 9, 7\\
DevKit & \texttt{Lab3/DevKit\_Lab3/ARDroneSimulinkDevKit\_V1.1/} (\texttt{start\_here\_CTRL.m}, \texttt{simulation/ARDroneTTSim.slx}, \texttt{wifiControl/ARDroneTT.slx}) & -\\
Lab 1 report & \texttt{labs-26-27/Lab1/report/} (your solutions, notably \texttt{q2.8.tex}, \texttt{q4.2.tex}) & -\\
\bottomrule
\end{tabularx}

\vspace{0.5ex}
There are no recorded video lectures in the repository.
The 25/26 notes say the Lab 3 devkit \texttt{guide.pdf} does not mention trajectory tracking at all, and that is right: the guide PDF in the devkit is the generic 2013 one, so the handout and this sheet are the only documentation.

% =====================================================================
\section*{Map: guide questions $\rightarrow$ areas}

\begin{tabularx}{\textwidth}{@{}l X@{}}
\toprule
Guide question & Area to read first\\
\midrule
2.1 (change of coordinates, $T$, $\theta_r$, $\phi_r$ from $u^*$) & Areas 0 and 1\\
2.2 (error system, input transformation, LQR gain) & Areas 1 and 2\\
2.3 (Lyapunov function $V_1$, $\dot V_1=-\tilde x^T Q^*\tilde x$) & Areas 2 and 3\\
2.4, 2.5 (disturbance, adaptation law) & Area 4\\
2.6 (global asymptotic stability of the augmented system) & Areas 3 and 4\\
2.7 (block diagram, $\bar K(s)$, link to PID) & Area 5\\
2.8 (yaw loop with zero steady-state error) & Area 5\\
3.1 (adapt the law to the DevKit inputs) & Areas 1 and 6\\
3.2, 3.3 (straight line, tuning) & Areas 6 and 7\\
3.4, 3.5, 3.6 (circle, yaw towards the centre) & Areas 6 and 7\\
3.7 (line-of-sight path following) & Area 8\\
\bottomrule
\end{tabularx}

% =====================================================================
\section*{Area 0 - Prerequisites (skim, about 45 min)}

\begin{tabularx}{\textwidth}{@{}C T X l@{}}
\toprule
\cb & Type & Source and topic & Pages\\
\midrule
\cb & \SLIDES & Ch1 deck: Euler angles and $R=R_z(\psi)R_y(\theta)R_x(\phi)$ (the convention used by \texttt{euler2R.m}, checked entry by entry) & skim\\
\cb & \SLIDES & Ch2 deck: quadrotor dynamic model & 15 to 18\\
\cb & \SLIDES & Ch2 deck: the quadrotor as a flat system (why $p_d,\dot p_d,\ddot p_d$ and $\psi_d$ are enough to define the motion) & 19 to 21\\
\cb & \LAB & \texttt{Lab1/report/q2.8.tex}: $G_y=\tfrac{g}{s^2}G_\phi$ and $G_x=-\tfrac{g}{s^2}G_\theta$, why the horizontal channels are harder, and why an outer loop must rotate by $\psi$ & whole file\\
\bottomrule
\end{tabularx}
\begin{selfcheck}
\chk{Which way is $z$ pointing in the course (NED), and what sign does a small pitch give for the $x$ acceleration?}
\chk{Why can a quadrotor only accelerate horizontally by tilting, and what does that imply for the cascade structure?}
\end{selfcheck}

% =====================================================================
\section*{Area 1 - Virtual input and cascade structure \hfill\normalfont\small (guide 2.1, 2.2, 3.1)}

\begin{tabularx}{\textwidth}{@{}C T X l@{}}
\toprule
\cb & Type & Source and topic & Pages\\
\midrule
\cb & \SLIDES & Ch3 deck: equations of motion (3 to 4), control objective (5), hierarchical or cascade control with the virtual input $u_T$, thrust magnitude and direction $r_{3d}$ (6 to 9) & 3 to 9\\
\cb & \SLIDES & Ch3 deck: position controller, error system and PD law (10 to 13) & 10 to 13\\
\cb & \SLIDES & Ch3 deck: attitude reference from $r_{3d}$ and $\psi$, via $[\cos\phi_d\sin\theta_d,\,-\sin\phi_d,\,\cos\phi_d\cos\theta_d]^T=R_z(-\psi)r_{3d}$ & 31 to 36\\
\cb & \SLIDES & Ch7-II deck: same decomposition for nonlinear tracking, $T_d=m\|u_T\|$, $r_{3d}=u_T/\|u_T\|$, and the full desired rotation matrix & 23 to 26 and 30 to 35\\
\cb & \NOTES & \texttt{02-control.pdf}: ``Control objective and hierarchical structure'' and ``Position controller design'' & Ch3 part\\
\bottomrule
\end{tabularx}
\begin{selfcheck}
\chk{Starting from $\ddot p=ge_3-\tfrac{T}{m}Re_3$, write $u^*$ so that $\ddot p=ge_3-R_z(\psi)u^*$, and invert it for $T,\theta_r,\phi_r$.}
\chk{Why does the yaw rotation $R_z(\psi)$ drop out of the roll and pitch references once you rotate by $R_z(-\psi)$?}
\chk{What does ``the inner loop is 5 to 10 times faster'' justify in the design, and where does the guide use it?}
\end{selfcheck}

% =====================================================================
\section*{Area 2 - LQR on the error system \hfill\normalfont\small (guide 2.2, 2.3, 3.3)}

\begin{tabularx}{\textwidth}{@{}C T X l@{}}
\toprule
\cb & Type & Source and topic & Pages\\
\midrule
\cb & \SLIDES & Ch4 deck: optimal control, Riccati equation, $K=R^{-1}B^TP$, cost $x^T(0)Px(0)$ & 26\\
\cb & \SLIDES & Ch4 deck: LQR design example, effect of the weights, Bryson's rule & 31 to 40\\
\cb & \BOOK & Chapter 6 ``Autopilot Design'', section 6.3 LQR Control & printed 115 to 117; PDF 123 to 125\\
\cb & \SLIDES & Ch4 deck: state-space, controllability, pole placement (so you can say what $(A,B)$ of the error system is) & 3 to 30\\
\cb & \NOTES & \texttt{02-control.pdf}: ``Optimal control: LQR'' and ``Tuning Q/R, Bryson's rule'' & Ch4 part\\
\bottomrule
\end{tabularx}
\begin{selfcheck}
\chk{Write $A,B$ of $\dot{\tilde x}=A\tilde x+Bu$ for the three axes, and say what the guide's cost $\int \tilde x^TQ\tilde x+\rho\|u\|^2$ means for $R$ in the slides.}
\chk{Show the closed-loop matrix $A-BK$ is Hurwitz without computing eigenvalues, using the theorem on slide 26.}
\chk{What does decreasing $\rho$ do to speed and to actuation, and which of the two limits you in the simulator?}
\end{selfcheck}

% =====================================================================
\section*{Area 3 - Lyapunov and LaSalle \hfill\normalfont\small (guide 2.3, 2.6)}

\begin{tabularx}{\textwidth}{@{}C T X l@{}}
\toprule
\cb & Type & Source and topic & Pages\\
\midrule
\cb & \SLIDES & Ch7 deck: stability definitions and Lyapunov's theorem & 14 to 22\\
\cb & \SLIDES & Ch7 deck: Krasovskii-LaSalle invariance principle, with the damped pendulum example & 23 to 31\\
\cb & \SLIDES & Ch7 deck: Lyapunov equation for linear systems, $A^TP+PA=-Q$ & 32 and 38 to 40\\
\cb & \SLIDES & Ch7-II deck: PD tracking proved GAS with LaSalle, and the choice of a different proportional gain & 27 to 29\\
\cb & \NOTES & \texttt{04-guidance-nonlinear.pdf}: ``Lyapunov's stability theorem'' and ``Krasovskii-LaSalle'' & Ch7 part I\\
\bottomrule
\end{tabularx}
\begin{selfcheck}
\chk{State what $\dot V\le 0$ gives, what $\dot V<0$ gives, and what you must add for the result to be global.}
\chk{When $\dot V$ is only negative semi-definite, which two things must you check to use LaSalle?}
\chk{Expand $\dot V_1$ for $V_1=\tilde x^TP\tilde x$ along $\dot{\tilde x}=(A-BK)\tilde x$ and use the Riccati equation. What is $Q^*$ in terms of $Q$, $P$, $B$, $\rho$?}
\end{selfcheck}

% =====================================================================
\section*{Area 4 - Adaptive disturbance rejection \hfill\normalfont\small (guide 2.4, 2.5, 2.6)}

\begin{tabularx}{\textwidth}{@{}C T X l@{}}
\toprule
\cb & Type & Source and topic & Pages\\
\midrule
\cb & \SLIDES & Ch7-II deck: Lyapunov analysis and adaptive control. The pattern is $V_2=V+\tfrac{1}{k_\theta}\tilde\theta^2$, then choose $\dot{\hat\theta}$ to cancel the cross term (14 to 15). Worked scalar example $\dot x=\theta x+u$ (16 to 18) & 14 to 18\\
\cb & \SLIDES & Ch7-II deck: backstepping, which rederives the PD law with a Lyapunov certificate (optional, context for why the law has the form it has) & 19 to 22\\
\cb & \NOTES & \texttt{04-guidance-nonlinear.pdf}: ``Lyapunov analysis and adaptive control'' & Ch7 part II\\
\cb & \NOTES & \texttt{labs-25-26/notes/03-lab3.pdf}: how the notes map 2.4 to 2.6 onto the pattern & Sec. 1\\
\bottomrule
\end{tabularx}
\begin{selfcheck}
\chk{Show that with $u^*=R_z(-\psi)(K\tilde x+ge_3+\hat w-\ddot p_d)$ the error dynamics become $\dot{\tilde x}=(A-BK)\tilde x+B\tilde w$.}
\chk{Why is $\dot{\tilde w}=-\dot{\hat w}$, and where is the assumption that $w$ is constant used?}
\chk{In the slides' scalar example the conclusion is only that $x\to 0$ and $\hat\theta$ is bounded (slides 17 and 18). Why can the lab system do better, i.e. what is different about the invariant set once $\tilde x\equiv 0$? Write the argument fully.}
\end{selfcheck}
Note on notation: the guide uses $k_w$ for the adaptation gain in 2.5 and again for the vertical proportional gain in 3.1.
They are two different numbers.
Name them differently in the report.

% =====================================================================
\section*{Area 5 - PID connection and the yaw loop \hfill\normalfont\small (guide 2.7, 2.8)}

\begin{tabularx}{\textwidth}{@{}C T X l@{}}
\toprule
\cb & Type & Source and topic & Pages\\
\midrule
\cb & \SLIDES & Ch3 deck: disturbance rejection with a PD loop (steady-state offset by the final value theorem) & 23 to 25\\
\cb & \SLIDES & Ch3 deck: integral action, from PD to PID, with experimental result & 26 to 29\\
\cb & \SLIDES & Ch3 deck: the unsolved exercise, wind as a velocity disturbance instead of an acceleration & 30\\
\cb & \SLIDES & Ch3 deck: state-space version with an integrator state, pole placement & 37 to 44\\
\cb & \NOTES & \texttt{02-control.pdf}: ``Disturbance rejection and integral action'' & Ch3 part\\
\bottomrule
\end{tabularx}
\begin{selfcheck}
\chk{With $\hat w$ driven by $\tilde x$, what is $\hat w(t)$ as an operator on the error? Use it to get $\bar K(s)$ and name the PID gains.}
\chk{Yaw model $\dot\psi=-a\psi+bu$ (not a pure integrator). Which controller gives zero steady-state error to a constant reference, and why does a pure P not?}
\end{selfcheck}

% =====================================================================
\section*{Area 6 - The DevKit: what the two empty blocks expect \hfill\normalfont\small (guide 3.1 to 3.7)}

This area has no slides.
It is a reading of \texttt{start\_here\_CTRL.m}, \texttt{setupTTSim.m}, \texttt{setupTT.m} and the two TT models.
Everything below was read from the files, except where marked \emph{not verified}.

\begin{tabularx}{\textwidth}{@{}C T X l@{}}
\toprule
\cb & Type & Source and topic & Where\\
\midrule
\cb & \DEVKIT & Run \texttt{start\_here\_CTRL.m}, answer \texttt{1} then \texttt{4}. It does \texttt{cd simulation; setupTTSim}, which sets \texttt{FMS.Ts=0.03}, \texttt{timeDelay=4*Ts} (0.12 s), \texttt{simDT=0.005}, loads the identified models and opens \texttt{ARDroneTTSim.slx} & scripts\\
\cb & \DEVKIT & Open \texttt{ARDroneTTSim.slx}. Blocks \texttt{MATLAB Function1} (= \texttt{desired\_trajectory}) and \texttt{Outer-loop Controller/MATLAB Function} (= \texttt{tt\_controller}) are the two to fill. Read the template code and comments in both & model\\
\cb & \DEVKIT & \texttt{lib/setupARModel.m}: the six identified inner-loop models (\texttt{ssRoll}, \texttt{ssRoll2V}, \texttt{ssPitch}, \texttt{ssPitch2U}, \texttt{ssYaw}, \texttt{ssH}). They are \texttt{.mat} objects: load them in MATLAB to see the poles and DC gains & lib\\
\cb & \LAB & \texttt{Lab1/report/q4.2.tex}: your identified closed-loop pitch model & Step 4\\
\bottomrule
\end{tabularx}

\subsection*{Interfaces of the two blocks}

\begin{tabularx}{\textwidth}{@{}l X@{}}
\toprule
Block & Signature (from the template code)\\
\midrule
\texttt{desired\_trajectory} & \texttt{[pd, dot\_pd, ddot\_pd, psi\_d] = f(on\_off, states, t)}. Persistent \texttt{t0}, \texttt{p0}, \texttt{previous\_status} are already there. \texttt{dt=t-t0} and \texttt{p0} are the two variables to build the trajectory from. The yaw is set to 0 when tracking is off.\\
\texttt{tt\_controller} & \texttt{[w\_r, dot\_psi\_r, theta\_r, phi\_r, dot\_xi] = f(states, pd, dot\_pd, ddot\_pd, psi\_d, xi)}. Already computed for you: \texttt{lambda}, \texttt{psi}, \texttt{R=euler2R(lambda)}, \texttt{Rz}, \texttt{psi\_e=normalize\_angle(psi-psi\_d)}, \texttt{p}, \texttt{dot\_p}.\\
\bottomrule
\end{tabularx}

\subsection*{Signal facts}

\begin{tabularx}{\textwidth}{@{}l X@{}}
\toprule
Item & What the models say\\
\midrule
\texttt{states} & \texttt{states(1:3)} Euler angles (rad), \texttt{states(4:6)} body velocity $(u,v,w)$, \texttt{states(7:9)} position. Position is \emph{integrated from velocity} by the block named ``Position estimation - inaccurate''. In simulation the drift is whatever the model gives. On the real drone it is the optical-flow warning in \texttt{setupTT.m}, and this is where Lab 2 matters.\\
Order of outputs & Top level wires the controller outputs to the drone block as: \texttt{w\_r} $\to$ input 5 (vertical speed), \texttt{dot\_psi\_r} $\to$ 4 (yaw rate), \texttt{theta\_r} $\to$ 3 (pitch), \texttt{phi\_r} $\to$ 2 (roll). Inputs 1, 6, 7 are fly, enable references, emergency stop.\\
\texttt{xi}, \texttt{dot\_xi} & The guide never mentions them. The controller block receives a 4-vector state \texttt{xi}, fed by a \emph{Discrete-Time Integrator} of \texttt{dot\_xi}: initial value \texttt{[0;0;0;0]}, sample time \texttt{FMS.Ts}, saturation $\pm 10$, and it is \textbf{reset while \texttt{on\_off}=0}. This is where the adaptive state $\hat w$ (3 components) goes, and plausibly one more integrator for the yaw loop of 2.8. The meaning of the fourth component is my inference from the size, not documented.\\
Controller selector & Inside \texttt{Outer-loop Controller}, \texttt{Manual Switch3} chooses between the old \emph{Baseline Controller} (hover cascade, with gains 1.5 for yaw, $-0.5$ forward, $0.4$ lateral) and the TT block. It is stored with \texttt{CurrentSetting=0}. \emph{Not verified} which input that is. Check on the diagram that the TT block is the one connected, with all outputs zero it will just sit there.\\
Inner-loop inputs & The drone block takes references in $[-1,1]$ and saturates them at $\pm 1$. The commands are \textbf{normalized, not radians}.\\
Real vs simulated & Both TT models contain identical copies of the two blocks. The Wi-Fi model only adds the Wi-Fi block (AT commands, packet decoding, yaw zeroed when references are enabled).\\
\bottomrule
\end{tabularx}

\begin{warnbox}
\textbf{1. Units of $\phi_r$, $\theta_r$, $w_r$, $\dot\psi_r$.}
The formulas of 2.1 give angles in radians, but the block feeds normalized commands.
Your own Lab 1 fit gives a DC gain of 0.414 from the pitch command to the pitch angle in rad (\texttt{q4.2.tex}), so a command of 1 is about 0.41 rad, and an angle reference in rad has to be divided by a factor of that order.
The factor for the simulation block itself is \emph{not verified}: read it from \texttt{dcgain(ssPitch)} and \texttt{dcgain(ssRoll)} after \texttt{setupARModel}.
The same question applies to $w_r$ and $\dot\psi_r$.
The 25/26 notes also flagged this and left it open.
\par\smallskip\textbf{2. Sign of the height.}
The template sets \texttt{p0 = [states(7:8); -1]}, so the third component of \texttt{pd} is $-1$ (NED $z$ for 1 m height).
But \texttt{tt\_controller} flips the sign of the measured position and velocity (\texttt{p(3)=-p(3)}, ``height is symmetric of z''), so its $p_3$ is the height, positive up.
Decide one convention for $p_{d,3}$ and check it with a first run, otherwise the vertical error can be off by 2 m.
\par\smallskip\textbf{3. On/off behaviour.}
\texttt{desired\_trajectory} resets \texttt{t0} and \texttt{p0} every time tracking is switched on, and holds \texttt{p0} with $\psi_d=0$ while it is off, so any extra state you add (for example for 3.7) must be reset in the same places.
\par\smallskip\textbf{4. Delay and sampling.}
The simulation includes a transport delay of $4T_s=0.12$ s on the drone side and the controller runs at 0.03 s.
Ch3 slides 14 to 22 give the tolerable delay as $\tau<PM/\omega_c$, which limits how aggressive $Q/\rho$ can be.
\end{warnbox}

\begin{selfcheck}
\chk{Which blocks are in the signal path from \texttt{tt\_controller} outputs to the drone, and what saturates the commands?}
\chk{Write the vector $p_d$ you would pass for hovering at 1 m, in the convention you chose, and what \texttt{p(3)} will be in the controller at that point.}
\end{selfcheck}

% =====================================================================
\section*{Area 7 - Reference trajectories \hfill\normalfont\small (guide 3.2, 3.4, 3.6)}

No new theory: this is calculus, plus decisions that the guide does not fix.

\begin{tabularx}{\textwidth}{@{}C T X l@{}}
\toprule
\cb & Type & Source and topic & Pages\\
\midrule
\cb & \SLIDES & Ch2 deck: flat outputs, why $\dot p_d,\ddot p_d$ are needed (feed-forward $\ddot p_d$ appears in the control law) & 19 to 21\\
\cb & \SLIDES & \texttt{ar23-ch10.pdf} (course wrap-up talk): last year's lab 3 example, circle path control with the vehicle heading towards the origin & 7\\
\bottomrule
\end{tabularx}
\begin{selfcheck}
\chk{Write $p_d,\dot p_d,\ddot p_d$ for the line (0.1 m/s) and for the circle ($R=1$ m, $\omega=0.1$ rad/s). What is the magnitude of $\ddot p_d$ on the circle, and why is it tiny?}
\chk{The circle must start at the drone's current position. Where is its centre? The guide leaves this free, so choose and justify.}
\chk{3.6: write $\psi_d$ so that the nose points at the centre, and its derivative. Does the position error equation in 2.4 depend on $\psi_d$ at all?}
\end{selfcheck}

% =====================================================================
\section*{Area 8 - Straight-line path following and LOS \hfill\normalfont\small (guide 3.7)}

\begin{tabularx}{\textwidth}{@{}C T X l@{}}
\toprule
\cb & Type & Source and topic & Pages\\
\midrule
\cb & \SLIDES & Ch8 deck: straight line description, path error $e_p$, cross-track error $e_{py}$ and its dynamics $\dot e_{py}=V_g\sin(\chi-\chi_q)$ & 5 to 7\\
\cb & \SLIDES & Ch8 deck: vector-field law $\chi^c=\chi_q-\chi^\infty\tfrac{2}{\pi}\arctan(k_{path}e_{py})$ and tuning of $k_{path}$ & 8 to 9\\
\cb & \SLIDES & Ch8 deck: Lyapunov proof $W=\tfrac12 e_{py}^2$ & 10 to 11\\
\cb & \SLIDES & Ch8 deck: lookahead distance, LOS and L1 (one slide, no derivation) & 19\\
\cb & \BOOK & Chapter 10, section 10.1 Straight-line Path Following & printed 198 to 203; PDF 206 to 211\\
\bottomrule
\end{tabularx}

\begin{warnbox}
The guide's formula $\psi_d=\arctan(-(y-y_d)/\Delta)$ is \textbf{not derived anywhere in the repository}.
I searched the book text: there is no line-of-sight path-following section in it (the ``line-of-sight'' hits are about vision and proportional navigation), and the slides give it one page.
You have to derive the geometry yourself, and the vector-field law above is the closest relative to compare with.
The standard external reference for LOS guidance is Fossen's work on guidance-based path following (for instance Breivik and Fossen); I did not check it here.
\end{warnbox}
\begin{selfcheck}
\chk{With $y_d$ the line's lateral position and $y$ the drone's, which sign of cross-track error makes $\psi_d$ positive, and does the drone turn towards the line?}
\chk{What do large and small $\Delta$ do to the intercept angle and to the speed of convergence? Compare with $k_{path}$ on slide 9.}
\chk{In 3.7, $\dot x_d=V\cos\psi_d$ and $\dot y_d=V\sin\psi_d$ are \emph{desired velocities}. How do they enter the $\dot p_d$ of the controller, and what is $\ddot p_d$ now?}
\end{selfcheck}

% =====================================================================
\section*{Area 9 - Closing synthesis (30 min)}

\begin{itemize}
\item \texttt{lectures-22-23/COURSE\_NOTES.pdf}: how modelling, control, estimation and guidance connect.
\item \texttt{labs-25-26/notes/03-lab3.pdf}, last two sections: the connections to the lectures and the open questions (some of them are answered in Area 6 above).
\item Lab 2 link: the TT controller uses a position and velocity that come from integrated optical-flow velocity. If you ever fly it, that is the Lab 2 estimation problem in practice.
\end{itemize}

% =====================================================================
\section*{Suggested schedule}

\begin{tabularx}{\textwidth}{@{}l X l@{}}
\toprule
Step & Area & Time\\
\midrule
1 & 0 and 1, then solve 2.1 and 2.2 on paper & about 2.5 h\\
2 & 2 and 3, then solve 2.3 & about 2 h\\
3 & 4, then solve 2.4 to 2.6 & about 2.5 h\\
4 & 5, then solve 2.7 and 2.8 & about 1.5 h\\
5 & 6, open the models, run the template once & about 1.5 h\\
6 & 7 and 8, derive the trajectories and the LOS geometry & about 1.5 h\\
7 & 9, then review all self-checks & about 1 h\\
\bottomrule
\end{tabularx}
The times are rough estimates, not measurements.
The theoretical questions are meant to be solved before the session, the simulations during it.

\subsection*{Minimal path if short on time}

\begin{enumerate}[nosep,leftmargin=1.5em]
\item Ch3 deck 3 to 13, Ch4 deck 26 and 31 to 40.
\item Ch7 deck 17 to 31, Ch7-II deck 14 to 18 and 27 to 29.
\item Area 6 of this sheet, with the three warnings.
\item Ch8 deck 5 to 11 and 19.
\end{enumerate}

\subsection*{Questions I would like to check you on later}

\begin{itemize}
\item Why is the whole position loop linear after the change of coordinates, and where does the non-linearity go?
\item Why is $\dot V_2$ only semi-definite, and what closes the proof in the lab system but not in the scalar example?
\item What does the adaptive state become in the Simulink block, and why is it reset when tracking is off?
\item Why does pointing the nose at the circle's centre not remove the tracking error, and what would?
\item What do $\Delta$ and $V$ do in 3.7?
\end{itemize}

\end{document}
